\chapter{Gallica, et ce qu'on a fait de germain-project}\label{ch:gallica}

Le second projet est celui dont un nouveau fonds du domaine public doit s'inspirer. Il a
repris la chaîne du projet Grothendieck sur une source dont les conditions sont presque
inverses : pas de droit d'auteur, un serveur ouvert, un protocole standard. Ce qui reste
difficile y apparaît d'autant mieux.

\section{La source}

La BnF a mis en ligne sur Gallica, entiers et gratuits, des manuscrits de mathématiciens :
papiers de travail, lettres, mises au net, catalogués par fonds et par numéro. Le projet
(\url{https://germain.commutator.io}, dépôt \texttt{germain-project}) en a dressé
l'inventaire à partir de la recherche de Gallica du 14 septembre 2026 : 78 volumes, 25\,258
vues, dix-neuf mathématiciens et deux albums de plusieurs mains, de Maurolico et Viète à
Libri et Prony. Les plus gros ensembles sont ceux de Fourier (20 volumes, 6\,173 vues), de
Mersenne (10 volumes, 3\,297 vues), de Libri (8 volumes, 2\,970 vues) et de Sophie Germain
(4 volumes, 1\,010 vues). Les détenteurs sont la BnF (Manuscrits et Arsenal) et, pour les
journaux de Prony, l'École des ponts.

Le site porte le nom de la première archive lue, celle de Sophie Germain.

\section{Les choix de départ}

\paragraph{Une seule règle de sélection.} Un volume n'entre au catalogue que si Gallica le
sert en entier. Un manuscrit non numérisé, ou que Gallica signale seulement chez un
partenaire, n'y figure pas, ni comme ligne ni comme vignette.

\paragraph{Les cahiers sont des regroupements du projet.} La BnF catalogue par fonds, pas
par mathématicien. Le site range les volumes en trois « cahiers » (\textsc{xvii}\ieme,
\textsc{xviii}\ieme, \textsc{xix}\ieme{} siècles), une section par mathématicien, et chaque
cahier dit que ce regroupement est le sien.

\paragraph{Aucune image stockée, nulle part.} Le volet du fac-similé demande à l'API IIIF
Image de Gallica la vue que le lecteur regarde, à la largeur du volet ; les octets vont de
la BnF au navigateur. Le \fichier{.gitignore} refuse \fichier{*.jpg} et \fichier{*.pdf}.
Faire défiler la transcription fait tourner les pages du fac-similé.

\paragraph{Pas de relais.} C'est le contraire de Montpellier, constaté le 10 septembre
2026 : les manifestes et les images IIIF répondent avec \texttt{Access-Control-Allow-Origin:
*}, le certificat est valide, et le champ \texttt{license} des manifestes renvoie aux
conditions de Gallica. Un script (\cmd{npm run check-sources}) vérifie à chaque fois que ces
conditions tiennent encore.

\section{Ce qui a coincé}

\paragraph{Le débit.} En septembre, une dizaine de requêtes en quelques secondes ont reçu
des coupures de connexion. Le 17 octobre 2026, le pare-feu placé devant la BnF a bloqué la
machine du projet sur toutes les adresses, avec une erreur~403, sans doute à cause des
copies de travail faites en septembre. Le script de copie a été réécrit en conséquence :
dix secondes entre deux images, avec une part aléatoire ; au plus 150 images par heure ;
arrêt complet au premier 403 ou à une erreur 429 qui persiste. Une adresse bloquée se
débloque en écrivant à la BnF, pas en réessayant.

\begin{warningblock}
C'est l'erreur la plus facile à commettre sur une source ouverte, et la plus coûteuse : une
copie de travail trop rapide fait bloquer le projet, et parfois toute une institution
derrière la même adresse. On règle la cadence avant le premier lot, pas après le blocage.
\end{warningblock}

\paragraph{La numérotation.} La plupart des vues de Gallica sont étiquetées « NP » (non
paginé). La transcription numérote donc par le numéro de vue de Gallica, qui pilote le
fac-similé, et note la foliotation de la bibliothèque seulement là où on la lit sur la
page, sans jamais la déduire.

\paragraph{La résolution.} Les images IIIF complètes sont grandes (environ $5200\times 7100$
pixels pour le manuscrit Français~9115). L'assistant les lit par bandes de 2\,400 pixels,
découpées dans la copie locale, ce qui ne coûte rien de plus à Gallica.

\paragraph{Les langues et les notations.} Français pour la plupart, latin pour Sluse,
Maurolico, Varignon et une partie de Mersenne et de Roberval. La transcription garde
l'orthographe, les abréviations et la notation de l'époque. La lecture modernisée traduit
la notation (les résidus de Legendre en congruences, par exemple) et signale chaque
traduction en note.

\section{Les droits, et ce qu'on a décidé}

Les manuscrits sont dans le domaine public : le dernier de ces mathématiciens est mort en
1869. Restent les conditions des numérisations. Celles de Gallica, lues le 23 septembre
2026 : la réutilisation non commerciale est libre, à condition de citer la source ; la
réutilisation commerciale est payante ; les métadonnées sont sous licence Etalab. Le site
est non commercial, sans publicité, et cite la source sous le volet et dans l'en-tête de
chaque fichier. Il soutient que ce qu'il extrait de la base de Gallica n'est pas
substantiel, et qu'il affiche Gallica parce que Gallica y invite.

Tout ce qui est au projet est sous CC0, transcriptions comprises. Pour une transcription
fidèle d'un texte du domaine public, CC0 dit seulement que le projet ne revendique rien :
une telle transcription vise l'exactitude et ne crée sans doute aucun droit nouveau. Pour
les lectures modernisées et le cahier d'exercices, qui sont des textes originaux, CC0
renonce à un droit qui existe.

Une règle encore : aucune transcription d'autrui n'est recopiée, ni prise comme point de
départ d'une lecture. Les procédures refusent de travailler sur la reproduction d'une page
dans une publication, et le site fabrique en retour du wikicode pour Wikisource.

\section{Ce qui a été produit}

D'après les fichiers, au début de novembre 2026 :
\begin{itemize}
  \item un peu plus de quatre-vingt-dix lots transcrits dans quinze volumes, chacun avec
        sa lecture modernisée. Les quatre volumes de Sophie Germain sont lus, dont le
        Français~9115 jusqu'à sa dernière vue ; des volumes de Mersenne, Viète, Descartes,
        Fermat, Roberval, Maurolico, Pascal et Lagrange sont commencés ;
  \item aucun lot n'est encore déclaré vérifié par une personne contre les vues ;
  \item 55 trouvailles, chacune écrite contre des sources nommées : 33 pistes, 19 déjà dans
        la littérature, 3 non encore cherchées ; la plupart portent sur les manuscrits
        eux-mêmes (36), puis sur les mathématiques (11) et l'histoire (8) ;
  \item environ 240 énoncés démontrés en Lean, tirés des volumes de Germain (arithmétique,
        courbure, élasticité), et un cahier de 42 exercices, lui aussi vérifié en Lean ;
  \item une frise des dix-neuf vies, une page des sources et des droits, une page de
        méthode, une page pour contribuer.
\end{itemize}

\begin{noteblock}
Le \fichier{README} du projet compte encore moins que les fichiers (seule Sophie Germain y
est donnée comme transcrite). Les compteurs du site, eux, sont tirés des fichiers. C'est
un défaut ordinaire et instructif : un chiffre écrit à la main dans un texte vieillit, un
chiffre calculé ne vieillit pas.
\end{noteblock}

\begin{exemple}
Pour un fonds du domaine public servi en IIIF, germain-project est le modèle à copier
presque tel quel : l'inventaire selon une règle unique, le fac-similé lu en direct, aucune
image dans le dépôt, un script de copie qui se nomme et ralentit, la vérification
périodique des en-têtes, et CC0 sur ce qui est à soi. Ses difficultés ne sont plus
juridiques ni techniques : elles tiennent à la cadence et à la relecture humaine.
\end{exemple}
